\section{Related Works}
\label{sec:formatting}

\paragraph{Video Diffusion Models}
Diffusion models~\cite{ho2020denoising, song2020denoising, song2020score, liu2022flow} have made great progress in image generation~\cite{dhariwal2021diffusion, nichol2021glide, ramesh2021zero, rombach2022high, saharia2022photorealistic, ramesh2022hierarchical, zhang2024creatilayout, zhang2025creatidesign, zhang2023adadiff, zhang2025blockdance, wang2025simplear, zhang2025reasongen}, which has led to the rapid development of video generation~\cite{ho2022video, wu2023tune, khachatryan2023text2video, guo2023animatediff, blattmann2023stable, zheng2024open, hong2022cogvideo, yang2024cogvideox, kong2024hunyuanvideo, xing2024survey, xing2024aid, xing2023vidiff, xie2025human2robot, du2024challenge, tu2024motioneditor, tu2024motionfollower, tu2025stableanimator, tu2025stableanimator++, tian2025unigen}.
VDM~\cite{ho2022video} is the first to apply diffusion models to video generation. 
Early works like AnimateDiff~\cite{guo2023animatediff} and SimDA~\cite{xing2024simda} attempt to insert temporal layers into pretrained T2I model for video generation.
Subsequently, VideoCrafter~\cite{chen2023videocrafter1} and SVD~\cite{blattmann2023stable} use large-scale and high-quality data for training, acheiving better performance.
However, these methods have trouble on generating long videos with high quality, mainly due to the inherent limitations of the UNet architecture.
The emergence of Sora~\cite{videoworldsimulators2024} is a significant success, demonstrating the potential of DiT~\cite{Peebles2022DiT} models to generate high-quality videos with tens of seconds.
Recent video generation methods~\cite{zheng2024open, yang2024cogvideox, kong2024hunyuanvideo} are mainly based on DiT architecture, and have achieved great success in the open-source community.
However, these methods rely solely on text or image guidance for video generation, lacking precise control over object or camera trajectory, which is crucial for high-quality video generation.

\vspace{-1em}
\paragraph{Trajectory Controllable Video Generation}
Trajectory Controllable Video Generation has recently garnered significant attention for its ability to precisely control object and camera trajectories during video synthesis.
Previous methods~\cite{yin2023dragnuwa, zhang2025tora, shi2024motion, li2025image}
integrate optical flow maps into video generation through a trajectory encoder.
% Points Trajectory
Recent works~\cite{wang2024motionctrl, zhou2025trackgo} suggest using point maps as a form of guidance.
MotionCtrl~\cite{wang2024motionctrl} processes point maps with a Gaussian filter and employs trainable encoders to encode object trajectories. Trackgo~\cite{zhou2025trackgo} represents objects using a few key points and injects this information via an encoder and a custom-designed adapter structure.
Other works~\cite{wang2024boximator,qiu2024freetraj, namekata2024sgi2v, jain2023peekaboo} employ bounding box to control object trajectories. Boximator~\cite{wang2024boximator} employs a trainable self-attention layer to fuse box and visual tokens inspired by GLIGEN~\cite{li2023gligen}. Some training-free methods~\cite{qiu2024freetraj, namekata2024sgi2v, jain2023peekaboo} purpose to modify attention layers or initial noised video latents to inject box signals.
Additionally, certain methods~\cite{wang2025cinemaster, gu2025das, fu20243dtrajmaster,wang2025levitor} explore the potential of 3D trajectories to achieve more sophisticated motion control. LeViTor~\cite{wang2025levitor} employs keypoint trajectory maps enriched with depth information, while others~\cite{wang2025cinemaster, gu2025das, fu20243dtrajmaster} construct custom 3D trajectories to represent object movements.
However, sparse trajectories lead to imprecise control on objects shape and size, while dense trajectories are difficult for users to provide.
In contrast, MagicMotion can control both dense and sparse trajectories, providing users with more flexible control over video generation.